\section{Covering and minimal-competitor lemmas}\label{app:transitions}
The two arguments below expand the topological and counting transitions
used in the finite-sheet and reconstruction theorems.

\begin{lemma}[A parameterized definable covering]\label{lem:param-cover}
Use the geometric hypotheses of Section~\ref{sec:reduction}, with $a(t)$
and $r(t)$ only continuous semialgebraic on $(0,t_0)$. Define
\[
 \mathcal B=\{(t,U):0<t<t_0,\ |U|<2r(t)\},\quad
 \mathcal X=\{(t,x):x\in\operatorname{int}N,
                  \ (t,F(x)-a(t))\in\mathcal B\}.
\]
The map $\Psi(t,x)=(t,F(x)-a(t))$ is a finite covering of $\mathcal B$,
proper over every compact subset of that base. Its fibre cardinality is
constant. Every covering component is the graph of a continuous
semialgebraic inverse sheet over $\mathcal B$, and those sheets exhaust
all inverse points. After shortening the time interval the sheets are
analytic for positive time. The centre arcs have convergent Puiseux
representatives after a finite power substitution in time. Every sheet
is defined on the closed smaller tube $|U|\le r(t)$.
\end{lemma}
\begin{proof}
There are no critical or boundary inverse points over $\mathcal B$.
The inverse function theorem for $F$, followed by the continuous shift
$a(t)$, makes $\Psi$ a local homeomorphism. This does not require
differentiating $a$ or $r$. If $K\subset\mathcal B$ is compact, its time
coordinates stay away from zero and $t_0$. A sequence in $\Psi^{-1}(K)$
has a subsequence converging in that compact time interval times $N$.
Its image stays in $K$; the limit cannot be on $\partial N$. Hence the
inverse image is compact. Each fibre is a compact discrete set and is
finite. A proper local homeomorphism is a covering: take disjoint local
inverse neighbourhoods of the finite fibre, and use compactness to rule
out additional inverse points over a shrinking target neighbourhood.
Its image is both open and closed in the connected base and is nonempty,
since the centre fibres are nonempty. It is therefore onto, with constant
finite cardinality.

The map $(t,U)\mapsto(t,U/r(t))$ identifies the base with an interval
times a ball; it is simply connected. Every covering component consequently
maps homeomorphically onto the base. The components are semialgebraic,
so their inverse graphs give definable labels. Continuous semialgebraic
one-variable arcs are analytic off a finite set. Shrinking the interval
removes that set, and the analytic inverse function theorem then makes
the sheets analytic at every positive-time point. The bounded centre
arcs have convergent Puiseux expansions by one-variable semialgebraic
preparation. Finally the closed smaller tube is contained in the open
larger tube, proving the last assertion.
\end{proof}

\begin{lemma}[Visibility of every minimal competitor]\label{lem:competitor}
Suppose a labelled profile $G$ has a fully visible $e$-endpoint
representation. Every representation of $G$ with $e$ endpoints is fully
visible. More generally, if $B'$ has independent columns, an active cone
with nonempty interior in $X^\circ$ contains a strict KKT point.
\end{lemma}
\begin{proof}
The intrinsic matrix difference between the largest and smallest local
quadratics has rank $e$. Every local difference for a competing
representation has range contained in $\ran B'$. Thus $B'$ has rank
$e$. There are at most $2^e$ active polynomials, while $G$ has exactly
$2^e$ distinct polynomials on open sets. Every one must therefore occur
on a full-dimensional active cone; a finite union of lower-dimensional
polyhedral cones cannot cover such an open set.

For a fixed active set $I$, each active coordinate is a row of
$-C_{II}^{-1}B_I^{\mathsf T}x$. No such row is identically zero, because
$B_I$ has independent columns. Each inactive multiplier is
$(B_j^{\mathsf T}-C_{jI}C_{II}^{-1}B_I^{\mathsf T})x$, a nonzero
functional because $B_j$ is not in the span of $B_I$. The finitely many
zero hyperplanes of these functionals cannot cover a full-dimensional
active cone in $X^\circ$. A point off them satisfies all its feasible
inequalities strictly. This proves both statements.
\end{proof}

\begin{lemma}[The ambient factor convention]\label{lem:distance-dimension}
For a positive block matrix with $e$ endpoints, take an invertible
$e\times e$ factor $L$ with $C=L^{\mathsf T}L$ and set
$T=-L^{-\mathsf T}B^{\mathsf T}\in\R^{e\times k}$. Then
\[
 G_Q(x)=x^{\mathsf T}(A-BC^{-1}B^{\mathsf T})x+
            \operatorname{dist}(Tx,L\R_+^e)^2.
\]
Two such invertible factors of the same $C$ differ by one left
orthogonal matrix, and that same matrix transports $T$. Positive
endpoint scales do not change the cone generated by the columns of $L$.
\end{lemma}
\begin{proof}
Expand $|Lz-Tx|^2$ and use $T^{\mathsf T}L=-B$. If
$L_1^{\mathsf T}L_1=L_2^{\mathsf T}L_2$, then
$O=L_2L_1^{-1}$ satisfies $O^{\mathsf T}O=I_e$ and $T_2=OT_1$.
Replacing $L$ by $LD$, for positive diagonal $D$, changes the lengths
but not the generating rays of its simplicial cone. These are statements
in the minimal ambient space $\R^e$, not assertions about unused
orthogonal directions in a larger factorization.
\end{proof}
